% Generated by simulations/python/experiments.py -- do not edit by hand
\begin{table}[htbp]\centering\small
\caption{Deadlock under unanimity as the assembly grows}\label{tab:deadlock_n}
\sbox{\tabbox}{\begin{tabular}{l>{\raggedleft\arraybackslash}p{0.140\linewidth}>{\raggedleft\arraybackslash}p{0.140\linewidth}>{\raggedleft\arraybackslash}p{0.131\linewidth}>{\raggedleft\arraybackslash}p{0.131\linewidth}>{\raggedleft\arraybackslash}p{0.131\linewidth}}
\toprule
Voters & Status quo Pareto-optimal (\%) & Unanimity: deadlock (\%) & QVC 1/n: fallback (\%) & QVC 25\%: fallback (\%) & QVC 33\%: fallback (\%) \\
\midrule
3 & 7.4 & 43.3 & 29.0 & 29.0 & 29.0 \\
5 & 26.7 & 67.4 & 51.7 & 14.5 & 14.5 \\
11 & 50.5 & 85.8 & 81.2 & 24.2 & 7.7 \\
21 & 70.1 & 94.6 & 97.4 & 15.5 & 6.4 \\
51 & 85.7 & 98.4 & 99.9 & 12.8 & 2.5 \\
101 & 92.8 & 99.6 & 100.0 & 8.9 & 0.8 \\
201 & 95.3 & 100.0 & 100.0 & 7.1 & 0.4 \\
\bottomrule
\end{tabular}}
\usebox{\tabbox}
\tabnotes{1,000 elections per row, spatial culture (2D), 8 options including the status quo. Status quo Pareto-optimal: the status quo lies in the convex hull of the voters' ideal points, so that no option whatsoever is preferred to it by every voter (Proposition~\ref{prop:gridlock}). Unanimity: no proposal is preferred to the status quo by every voter, which stays. QVC: share of the elections in which every option is vetoed by the quota and the least-vetoed options are kept; QVC still decides. Monte Carlo standard errors at most 1.6 points.}
{Monte Carlo simulation, cultures of \texttt{cultures.py}, rules of \texttt{voting.py}; \texttt{experiments.py}.}
\end{table}
